\documentclass[aps,preprint]{revtex4}%
\usepackage{amsfonts}%
\usepackage{amsmath}%
\setcounter{MaxMatrixCols}{30}%
\usepackage{amssymb}%
\usepackage{graphicx}
%TCIDATA{OutputFilter=latex2.dll}
%TCIDATA{Version=4.10.0.2363}
%TCIDATA{CSTFile=revtex4.cst}
%TCIDATA{Created=Tuesday, June 17, 2003 17:08:15}
%TCIDATA{LastRevised=Tuesday, June 17, 2003 17:15:20}
%TCIDATA{<META NAME="GraphicsSave" CONTENT="32">}
%TCIDATA{<META NAME="DocumentShell" CONTENT="Articles\SW\REVTeX 4">}
\newtheorem{theorem}{Theorem}
\newtheorem{acknowledgement}[theorem]{Acknowledgement}
\newtheorem{algorithm}[theorem]{Algorithm}
\newtheorem{axiom}[theorem]{Axiom}
\newtheorem{claim}[theorem]{Claim}
\newtheorem{conclusion}[theorem]{Conclusion}
\newtheorem{condition}[theorem]{Condition}
\newtheorem{conjecture}[theorem]{Conjecture}
\newtheorem{corollary}[theorem]{Corollary}
\newtheorem{criterion}[theorem]{Criterion}
\newtheorem{definition}[theorem]{Definition}
\newtheorem{example}[theorem]{Example}
\newtheorem{exercise}[theorem]{Exercise}
\newtheorem{lemma}[theorem]{Lemma}
\newtheorem{notation}[theorem]{Notation}
\newtheorem{problem}[theorem]{Problem}
\newtheorem{proposition}[theorem]{Proposition}
\newtheorem{remark}[theorem]{Remark}
\newtheorem{solution}[theorem]{Solution}
\newtheorem{summary}[theorem]{Summary}
\newenvironment{proof}[1][Proof]{\noindent\textbf{#1.} }{\ \rule{0.5em}{0.5em}}
\begin{document}
\preprint{ }
\title[Canonical Form]{Canonical Form of Complex Antisymmetric Matrices}
\author{Brian Weiner}
\affiliation{PSU}

\begin{abstract}


\end{abstract}
\date{06/17/2003}
\startpage{1}
\maketitle
\tableofcontents


\section{Antisymmetric Operators}

\subsection{Real Antisymmetric Matrices}

Let $\mathcal{A}$ be a \emph{real} compact antisymmetric operator
$\in\mathcal{B}_{0}\left(  V\right)  $, where $V$ is a Hilbert space then it
has the canonical expansion
\begin{equation}
\mathcal{A=}\sum_{1\leq j\leq s}\alpha_{j}\left\{  \left\vert s_{j}%
\right\rangle \left\langle s_{j+s}\right\vert -\left\vert s_{_{j+s}%
}\right\rangle \left\langle s_{j}\right\vert \right\} \label{anti}%
\end{equation}
As the operator is real $\mathcal{A=A}^{\ast}$, (see Appendix \ref{star} for
the definition of complex conjugation *), it can be converted into a compact
self adjoint purely imaginary operator, $\mathcal{H}$, with the properties
\begin{align}
\mathcal{H}  & =i\mathcal{A}\nonumber\\
\mathcal{H}^{\dagger}  & =-i\mathcal{A}^{t}=i\mathcal{A=H}\nonumber\\
\mathcal{H}^{\ast}  & =-i\mathcal{A}=i\mathcal{A}^{t}\,\mathcal{=H}^{t}%
\end{align}
and
\begin{equation}
\mathcal{H}=i\mathcal{A=}\sum_{1\leq j\leq s}\alpha_{j}\left\{  \left\vert
v_{j}\right\rangle \left\langle v_{j}\right\vert -\left\vert v_{j}^{\ast
}\right\rangle \left\langle v_{j}^{\ast}\right\vert \right\}  ;\qquad
\alpha_{j}\in\mathbb{R}\label{sa}%
\end{equation}
where $\left\{  \left\vert v_{j}\right\rangle ,\left\vert v_{j}^{\ast
}\right\rangle ;1\leq j\leq\infty\right\}  $ is a conb for $\mathcal{H}$. Eq.
$\left(  \ref{sa}\right)  $ is seen by observing that
\begin{equation}
\mathcal{H}\left\vert v\right\rangle =\alpha\left\vert v\right\rangle
\Rightarrow\mathcal{H}^{\ast}\left\vert v^{\ast}\right\rangle =\alpha
\left\vert v^{\ast}\right\rangle \Rightarrow\mathcal{H}^{t}\left\vert v^{\ast
}\right\rangle =\alpha\left\vert v^{\ast}\right\rangle \Rightarrow
\mathcal{H}\left\vert v^{\ast}\right\rangle =-\alpha\left\vert v^{\ast
}\right\rangle
\end{equation}
It is important to note that the orthonormal eigenvectors of $\mathcal{H}$ are
only defined up to an overall phase factors $e^{i\varphi_{j}}$ which \emph{do
not} preserve the complex conjugate relationship between eigenvectors
associated with paired eigenvalues of opposite sign. These different bases are
related to each other by
\begin{equation}
\left(  \left\vert v_{1}^{\prime}\right\rangle \cdots\left\vert v_{2s}%
^{\prime}\right\rangle \right)  =\left(  \left\vert v_{1}\right\rangle
\cdots\left\vert v_{2s}\right\rangle \right)  \left(
\begin{array}
[c]{ccc}%
e^{i\varphi_{1}} & \mathbf{0} & \mathbf{0}\\
\mathbf{0} & \ddots & \mathbf{0}\\
\mathbf{0} & \mathbf{0} & e^{i\varphi_{2s}}%
\end{array}
\right)
\end{equation}
Using a basis with complex conjugate symmetry we obtain
\begin{equation}
\mathcal{A=-}\sum_{1\leq j\leq s}i\alpha_{j}\left\{  \left\vert v_{j}%
\right\rangle \left\langle v_{j}\right\vert -\left\vert v_{j}^{\ast
}\right\rangle \left\langle v_{j}^{\ast}\right\vert \right\}  ;\qquad
\alpha_{i}\in\mathbb{R}%
\end{equation}
We can then define a \emph{real} conb $\left\{  \left\vert s_{j}\right\rangle
;1\leq j\leq\infty\right\}  $ i.e. $\left\vert s_{j}\right\rangle =\left\vert
s_{j}^{\ast}\right\rangle $ as
\begin{align}
\left\vert s_{j}\right\rangle  & =\frac{1}{\sqrt{2}}\left(  \left\vert
v_{j}\right\rangle +\left\vert v_{j}^{\ast}\right\rangle \right)  =\sqrt
{2}\operatorname{Re}\left\vert v_{j}\right\rangle \nonumber\\
\left\vert s_{j+s}\right\rangle  & =\frac{-i}{\sqrt{2}}\left(  \left(
\left\vert v_{j}\right\rangle -\left\vert v_{j}^{\ast}\right\rangle \right)
\right)  =\sqrt{2}\operatorname{Im}\left\vert v_{j}\right\rangle
\end{align}
so that
\begin{align}
\mathcal{A}\left\{  \frac{1}{\sqrt{2}}\left(  \left\vert v_{j}\right\rangle
+\left\vert v_{j}^{\ast}\right\rangle \right)  \right\}   & =-i\alpha_{j}%
\frac{1}{\sqrt{2}}\left(  \left\vert v_{j}\right\rangle -\left\vert
v_{j}^{\ast}\right\rangle \right) \nonumber\\
& =\alpha_{j}\left\{  \frac{-i}{\sqrt{2}}\left(  \left(  \left\vert
v_{j}\right\rangle -\left\vert v_{j}^{\ast}\right\rangle \right)  \right)
\right\}
\end{align}
and
\begin{align}
\mathcal{A}\left\{  \frac{-i}{\sqrt{2}}\left(  \left(  \left\vert
v_{j}\right\rangle -\left\vert v_{j}^{\ast}\right\rangle \right)  \right)
\right\}   & =\frac{-\alpha_{j}}{\sqrt{2}}\left\vert v_{j}\right\rangle
-\frac{\alpha_{j}}{\sqrt{2}}\left\vert v_{j}^{\ast}\right\rangle \nonumber\\
& =-\alpha_{j}\left\{  \frac{1}{\sqrt{2}}\left(  \left\vert v_{j}\right\rangle
+\left\vert v_{j}^{\ast}\right\rangle \right)  \right\}
\end{align}
thus
\begin{equation}
\mathcal{A=}\sum_{1\leq j\leq s}\alpha_{j}\left\{  \left\vert s_{j}%
\right\rangle \left\langle s_{j+s}\right\vert -\left\vert s_{_{j+s}%
}\right\rangle \left\langle s_{j}\right\vert \right\}
\end{equation}
which justifies the expansion in eq. $\left(  \ref{anti}\right)  $. Using the
bases $\left\{  \left\vert s_{j}\right\rangle ;1\leq j\leq\infty\right\}  $ we
can produce the canonical matrix representation of $\mathcal{A}$
\begin{equation}
\mathcal{A}\left(  \left\vert \mathbf{s}\right\rangle \right)  =\left\vert
\mathbf{s}\right\rangle \left(
\begin{array}
[c]{cc}%
\mathbb{O} & \mathbf{\alpha}\\
-\mathbf{\alpha} & \mathbb{O}%
\end{array}
\right)
\end{equation}
where
\begin{equation}
\mathbf{\alpha=}\left(
\begin{array}
[c]{ccc}%
\alpha_{1} & \mathbb{O} & \mathbb{O}\\
\mathbb{O} & \ddots & \mathbb{O}\\
\mathbb{O} & \mathbb{O} & \alpha_{s}%
\end{array}
\right)
\end{equation}
One can reorder this basis to obtain
\begin{equation}
\mathcal{A}\left(  \left\vert s_{1}\right\rangle \left\vert s_{1+s}%
\right\rangle \cdots\left\vert s_{s}\right\rangle \left\vert s_{s+s}%
\right\rangle \right)  =\left(  \left\vert s_{1}\right\rangle \left\vert
s_{1+s}\right\rangle \cdots\left\vert s_{s}\right\rangle \left\vert
s_{s+s}\right\rangle \right)  \left(
\begin{array}
[c]{ccccc}%
0 & \alpha_{1} & \mathbb{O} & \mathbb{O} & \mathbb{O}\\
-\alpha_{1} & 0 & \mathbb{O} & \mathbb{O} & \mathbb{O}\\
\mathbb{O} & \mathbb{O} & \ddots & \mathbb{O} & \mathbb{O}\\
\mathbb{O} & \mathbb{O} & \mathbb{O} & 0 & \alpha_{s}\\
\mathbb{O} & \mathbb{O} & \mathbb{O} & -\alpha_{s} & 0
\end{array}
\right)
\end{equation}
and the expansion
\begin{equation}
\mathcal{A}=\left(  \left\vert s_{1}\right\rangle \left\vert s_{1+s}%
\right\rangle \cdots\left\vert s_{s}\right\rangle \left\vert s_{s+s}%
\right\rangle \right)  \left(
\begin{array}
[c]{ccccc}%
0 & \alpha_{1} & \mathbb{O} & \mathbb{O} & \mathbb{O}\\
-\alpha_{1} & 0 & \mathbb{O} & \mathbb{O} & \mathbb{O}\\
\mathbb{O} & \mathbb{O} & \ddots & \mathbb{O} & \mathbb{O}\\
\mathbb{O} & \mathbb{O} & \mathbb{O} & 0 & \alpha_{s}\\
\mathbb{O} & \mathbb{O} & \mathbb{O} & -\alpha_{s} & 0
\end{array}
\right)  \left(
\begin{array}
[c]{c}%
\left\vert s_{1}\right\rangle \\
\left\vert s_{1+s}\right\rangle \\
\vdots\\
\left\vert s_{s}\right\rangle \\
\left\vert s_{s+s}\right\rangle
\end{array}
\right) \label{reex}%
\end{equation}
The row vector of basis vectors, $\left(  \left\vert s_{1}\right\rangle
\left\vert s_{1+s}\right\rangle \cdots\left\vert s_{s}\right\rangle \left\vert
s_{s+s}\right\rangle \right)  $, can be viewed as a real orthogonal matrix
whose columns are the real basis vectors $\left\{  \cdots,\left\vert
s_{j}\right\rangle ,\left\vert s_{s+j}\right\rangle ,\cdots\right\}  $.

The complex basis $\left\{  \left\vert v_{j}\right\rangle ,\left\vert
v_{j}^{\ast}\right\rangle ;1\leq j\leq\infty\right\}  $ diagonalizes
$\mathcal{A}$ to give
\begin{equation}
\mathcal{A}\left(  \left\vert \mathbf{v}\right\rangle ,\left\vert
\mathbf{v}^{\ast}\right\rangle \right)  =\left(  \left\vert \mathbf{v}%
\right\rangle ,\left\vert \mathbf{v}^{\ast}\right\rangle \right)  \left(
\begin{array}
[c]{cc}%
-i\mathbf{\alpha} & \mathbb{O}\\
\mathbb{O} & i\mathbf{\alpha}%
\end{array}
\right)
\end{equation}
this basis can also be reordered to produce
\begin{equation}
\mathcal{A}\left(  \left\vert \left\vert v_{1}\right\rangle \left\vert
v_{1}^{\ast}\right\rangle \cdots\left\vert v_{s}\right\rangle \left\vert
v_{s}^{\ast}\right\rangle \right\rangle \right)  =\left(  \left\vert
v_{1}\right\rangle \left\vert v_{1}^{\ast}\right\rangle \cdots\left\vert
v_{s}\right\rangle \left\vert v_{s}^{\ast}\right\rangle \right)  \left(
\begin{array}
[c]{ccccc}%
-i\alpha_{1} & 0 & \mathbb{O} & \mathbb{O} & \mathbb{O}\\
0 & i\alpha_{1} & \mathbb{O} & \mathbb{O} & \mathbb{O}\\
\mathbb{O} & \mathbb{O} & \ddots & \mathbb{O} & \mathbb{O}\\
\mathbb{O} & \mathbb{O} & \mathbb{O} & -i\alpha_{s} & 0\\
\mathbb{O} & \mathbb{O} & \mathbb{O} & 0 & i\alpha_{s}%
\end{array}
\right)
\end{equation}
and the expansion%
\begin{equation}
\mathcal{A}=\left(  \left\vert v_{1}\right\rangle \left\vert v_{1}^{\ast
}\right\rangle \cdots\left\vert v_{s}\right\rangle \left\vert v_{s}^{\ast
}\right\rangle \right)  \left(
\begin{array}
[c]{ccccc}%
-i\alpha_{1} & 0 & \mathbb{O} & \mathbb{O} & \mathbb{O}\\
0 & i\alpha_{1} & \mathbb{O} & \mathbb{O} & \mathbb{O}\\
\mathbb{O} & \mathbb{O} & \ddots & \mathbb{O} & \mathbb{O}\\
\mathbb{O} & \mathbb{O} & \mathbb{O} & -i\alpha_{s} & 0\\
\mathbb{O} & \mathbb{O} & \mathbb{O} & 0 & i\alpha_{s}%
\end{array}
\right)  \left(
\begin{array}
[c]{c}%
\left\vert v_{1}\right\rangle \\
\left\vert v_{1}^{\ast}\right\rangle \\
\vdots\\
\left\vert v_{s}\right\rangle \\
\left\vert v_{s}^{\ast}\right\rangle
\end{array}
\right) \label{ccex}%
\end{equation}
The row vector of basis vectors, $\left(  \left\vert v_{1}\right\rangle
\left\vert v_{1}^{\ast}\right\rangle \cdots\left\vert v_{s}\right\rangle
\left\vert v_{s}^{\ast}\right\rangle \right)  $, can be viewed as an unitary
matrix whose columns are the basis vectors $\left\{  \cdots,\left\vert
v_{j}\right\rangle ,\left\vert v_{j}^{\ast}\right\rangle ,\cdots\right\}  $.
Note that the expansions of eqs $\left(  \ref{reex}\right)  $ and $\left(
\ref{ccex}\right)  $ both give $\mathcal{A}$ thus in eq $\left(
\ref{reex}\right)  $ all the imaginary terms cancel.

\subsection{Complex Antisymmetric Matrices}

If the \emph{self adjoint} real and complex parts of a complex matrix commute,
(note that these parts \emph{need not }be real), then the matrix defines a
representation of a normal operator. If such a matrix is antisymmetric it can
be brought to complex canonical form by a \emph{real orthogonal}
transformation as above.

This can be seen from the following: Let
\begin{equation}
\left[  \mathcal{A},\mathcal{A}^{\dagger}\right]  =0
\end{equation}
then the \emph{self adjoint }real and imaginary parts of $\mathcal{A}$ are
defined as
\begin{subequations}
\begin{align}
\operatorname{Re}\mathcal{A}  & =\frac{1}{2}\left(  \mathcal{A}+\mathcal{A}%
^{\dagger}\right) \\
\operatorname{Im}\mathcal{A}  & =\frac{-i}{2}\left(  \mathcal{A}%
-\mathcal{A}^{\dagger}\right) \\
\mathcal{A}  & =\operatorname{Re}\mathcal{A}+i\operatorname{Im}\mathcal{A}%
\end{align}
which commute i.e.
\end{subequations}
\begin{equation}
\left[  \operatorname{Re}\mathcal{A},\operatorname{Im}\mathcal{A}\right]  =0
\end{equation}
If $\mathcal{A}$ is antisymmetric
\begin{equation}
\mathcal{A}=-\mathcal{A}^{t}\label{at}%
\end{equation}
(the transpose operation is defined in terms of the adjoint and complex
conjugation operations see Appendix \ref{star})\ then
\begin{align}
\operatorname{Re}\mathcal{A\,}  & =-\operatorname{Re}\mathcal{A}%
^{t}\nonumber\\
\operatorname{Im}\mathcal{A\,}  & =-\operatorname{Im}\mathcal{A}^{t}.
\end{align}
and
\begin{align}
\operatorname{Re}\mathcal{A}  & =\operatorname{Re}\mathcal{A}^{\dagger
}=\operatorname{Re}\mathcal{A}^{t\ast}=-\left(  \operatorname{Re}%
\mathcal{A}\right)  ^{\ast}\nonumber\\
\operatorname{Im}\mathcal{A}  & =\operatorname{Im}\mathcal{A}^{\dagger
}=\operatorname{Im}\mathcal{A}^{t\ast}=-\left(  \operatorname{Im}%
\mathcal{A}\right)  ^{\ast}%
\end{align}
Hence $\operatorname{Re}\mathcal{A}$ and $\operatorname{Im}\mathcal{A}$ are
purely imaginary i.e. products of $i$ and real antisymmetric matrices i.e.
\begin{align}
\operatorname{Re}\mathcal{A}  & =i\mathcal{M}\nonumber\\
\operatorname{Im}\mathcal{A}  & =-i\mathcal{N}%
\end{align}
where
\begin{align}
\mathcal{M}  & =-\mathcal{M}^{t}=\mathcal{M}^{\ast}\nonumber\\
\mathcal{N}  & =-\mathcal{N}^{t}=\mathcal{N}^{\ast}%
\end{align}
so that
\begin{equation}
\mathcal{A}=i\mathcal{M}+\mathcal{N}=\mathcal{N}+i\mathcal{M}%
\end{equation}
and
\begin{align}
\mathcal{N}  & =\frac{1}{2}\left(  \mathcal{A}-\mathcal{A}^{\dagger}\right)
=\frac{1}{2}\left(  \mathcal{A}+\mathcal{A}^{\ast}\right)  =-\mathcal{N}%
^{\dagger}=\mathcal{N}^{\ast}\nonumber\\
\mathcal{M}  & =\frac{-i}{2}\left(  \mathcal{A}+\mathcal{A}^{\dagger}\right)
=\frac{-i}{2}\left(  \mathcal{A}-\mathcal{A}^{\ast}\right)  =-\mathcal{M}%
^{\dagger}=\mathcal{M}^{\ast}%
\end{align}
confirming that $\mathcal{M}$ and $\mathcal{N}$ are real. (Note that the real
part $\mathcal{N}$ of $\mathcal{A}$ is related to self adjoint
\emph{imaginary} part of $\mathcal{A}$ while the imaginary part $\mathcal{M}$
of $\mathcal{A}$ is related to self adjoint \emph{real} part of $\mathcal{A}%
$!). The matrices $\mathcal{M}$ and $\mathcal{N}$ commute i.e
\begin{equation}
\left[  \mathcal{M},\mathcal{N}\right]  =0
\end{equation}
as do the self adjoint real and complex parts of $\mathcal{A}.$ $\lceil$The
above shows that As $\mathcal{A=-A}^{t}$ it it follows that $\left[
\mathcal{A},\mathcal{A}^{\ast}\right]  =0$ which then implies that
$\mathcal{A}$ is normal$\rfloor.$

$\operatorname{Re}\mathcal{A}$ and $\operatorname{Im}\mathcal{A}$ can be
simultaneously diagonalized to give
\begin{align}
\operatorname{Re}\mathcal{A}  & =i\mathcal{M}=\sum_{1\leq j\leq s}\mu
_{j}\left\{  \left\vert v_{j}\right\rangle \left\langle v_{j}\right\vert
-\left\vert v_{j}^{\ast}\right\rangle \left\langle v_{j}^{\ast}\right\vert
\right\}  ;\qquad\mu_{j}\in\mathbb{R}\nonumber\\
& \Rightarrow\mathcal{M}=\mathcal{-}\sum_{1\leq j\leq s}i\mu_{j}\left\{
\left\vert v_{j}\right\rangle \left\langle v_{j}\right\vert -\left\vert
v_{j}^{\ast}\right\rangle \left\langle v_{j}^{\ast}\right\vert \right\}
\end{align}
and
\begin{align}
\operatorname{Im}\mathcal{A}  & =-i\mathcal{N}=\sum_{1\leq j\leq s}%
\upsilon_{j}\left\{  \left\vert v_{j}\right\rangle \left\langle v_{j}%
\right\vert -\left\vert v_{j}^{\ast}\right\rangle \left\langle v_{j}^{\ast
}\right\vert \right\}  ;\qquad\upsilon_{j}\in\mathbb{R}\nonumber\\
& \Rightarrow\mathcal{N}=\sum_{1\leq j\leq s}i\upsilon_{j}\left\{  \left\vert
v_{j}\right\rangle \left\langle v_{j}\right\vert -\left\vert v_{j}^{\ast
}\right\rangle \left\langle v_{j}^{\ast}\right\vert \right\}
\end{align}
Then we can use the theory for real antisymmetric matrices to obtain
\begin{equation}
\mathcal{M}=\sum_{1\leq j\leq s}\mu_{j}\left\{  \left\vert s_{j}\right\rangle
\left\langle s_{j+s}\right\vert -\left\vert s_{j}\right\rangle \left\langle
s_{j+s}\right\vert \right\}
\end{equation}
and
\begin{equation}
\mathcal{N}=\sum_{1\leq j\leq s}\left(  -\upsilon_{j}\right)  \left\{
\left\vert s_{j}\right\rangle \left\langle s_{j+s}\right\vert -\left\vert
s_{_{j+s}}\right\rangle \left\langle s_{j}\right\vert \right\}
\end{equation}
which gives using
\begin{equation}
\mathcal{A}=\operatorname{Re}\mathcal{A}+i\operatorname{Im}\mathcal{A}%
=\mathcal{N}+i\mathcal{M}%
\end{equation}
that
\begin{align}
\mathcal{A}  & =\sum_{1\leq j\leq s}\left(  -\upsilon_{j}+i\mu_{j}\right)
\left\{  \left\vert s_{j}\right\rangle \left\langle s_{j+s}\right\vert
-\left\vert s_{_{j+s}}\right\rangle \left\langle s_{j}\right\vert \right\}
\nonumber\\
& \equiv\sum_{1\leq j\leq s}\eta_{j}\left\{  \left\vert s_{j}\right\rangle
\left\langle s_{j+s}\right\vert -\left\vert s_{_{j+s}}\right\rangle
\left\langle s_{j}\right\vert \right\}
\end{align}
where $\eta_{j}=-\upsilon_{j}+i\mu_{j}$. This gives the matrix representation
\begin{equation}
\mathcal{A}\left(  \left\vert s_{1}\right\rangle \left\vert s_{1+s}%
\right\rangle \cdots\left\vert s_{s}\right\rangle \left\vert s_{s+s}%
\right\rangle \right)  =\left(  \left\vert s_{1}\right\rangle \left\vert
s_{1+s}\right\rangle \cdots\left\vert s_{s}\right\rangle \left\vert
s_{s+s}\right\rangle \right)  \left(
\begin{array}
[c]{ccccc}%
0 & \eta_{1} & \mathbb{O} & \mathbb{O} & \mathbb{O}\\
-\eta_{1} & 0 & \mathbb{O} & \mathbb{O} & \mathbb{O}\\
\mathbb{O} & \mathbb{O} & \ddots & \mathbb{O} & \mathbb{O}\\
\mathbb{O} & \mathbb{O} & \mathbb{O} & 0 & \eta_{s}\\
\mathbb{O} & \mathbb{O} & \mathbb{O} & -\eta_{s} & 0
\end{array}
\right)
\end{equation}
and the expansion
\begin{equation}
\mathcal{A}=\left(  \left\vert s_{1}\right\rangle \left\vert s_{1+s}%
\right\rangle \cdots\left\vert s_{s}\right\rangle \left\vert s_{s+s}%
\right\rangle \right)  \left(
\begin{array}
[c]{ccccc}%
0 & \eta_{1} & \mathbb{O} & \mathbb{O} & \mathbb{O}\\
-\eta_{1} & 0 & \mathbb{O} & \mathbb{O} & \mathbb{O}\\
\mathbb{O} & \mathbb{O} & \ddots & \mathbb{O} & \mathbb{O}\\
\mathbb{O} & \mathbb{O} & \mathbb{O} & 0 & \eta_{s}\\
\mathbb{O} & \mathbb{O} & \mathbb{O} & -\eta_{s} & 0
\end{array}
\right)  \left(
\begin{array}
[c]{c}%
\left\vert s_{1}\right\rangle \\
\left\vert s_{1+s}\right\rangle \\
\vdots\\
\left\vert s_{s}\right\rangle \\
\left\vert s_{s+s}\right\rangle
\end{array}
\right)
\end{equation}
We can also observe that
\begin{equation}
\mathcal{A}=-i\sum_{1\leq j\leq s}\eta_{j}\left\{  \left\vert v_{j}%
\right\rangle \left\langle v_{j}\right\vert -\left\vert v_{j}^{\ast
}\right\rangle \left\langle v_{j}^{\ast}\right\vert \right\}
\end{equation}
This gives the matrix representation
\begin{equation}
\mathcal{A}\left(  \left\vert \left\vert v_{1}\right\rangle \left\vert
v_{1}^{\ast}\right\rangle \cdots\left\vert v_{s}\right\rangle \left\vert
v_{s}^{\ast}\right\rangle \right\rangle \right)  =\left(  \left\vert
\left\vert v_{1}\right\rangle \left\vert v_{1}^{\ast}\right\rangle
\cdots\left\vert v_{s}\right\rangle \left\vert v_{s}^{\ast}\right\rangle
\right\rangle \right)  \left(
\begin{array}
[c]{ccccc}%
-i\eta_{1} & 0 & \mathbb{O} & \mathbb{O} & \mathbb{O}\\
0 & i\eta_{1} & \mathbb{O} & \mathbb{O} & \mathbb{O}\\
\mathbb{O} & \mathbb{O} & \ddots & \mathbb{O} & \mathbb{O}\\
\mathbb{O} & \mathbb{O} & \mathbb{O} & -i\eta_{s} & 0\\
\mathbb{O} & \mathbb{O} & \mathbb{O} & 0 & i\eta_{s}%
\end{array}
\right)
\end{equation}
and the expansion%
\begin{equation}
\mathcal{A}=\left(  \left\vert \left\vert v_{1}\right\rangle \left\vert
v_{1}^{\ast}\right\rangle \cdots\left\vert v_{s}\right\rangle \left\vert
v_{s}^{\ast}\right\rangle \right\rangle \right)  \left(
\begin{array}
[c]{ccccc}%
-i\eta_{1} & 0 & \mathbb{O} & \mathbb{O} & \mathbb{O}\\
0 & i\eta_{1} & \mathbb{O} & \mathbb{O} & \mathbb{O}\\
\mathbb{O} & \mathbb{O} & \ddots & \mathbb{O} & \mathbb{O}\\
\mathbb{O} & \mathbb{O} & \mathbb{O} & -i\eta_{s} & 0\\
\mathbb{O} & \mathbb{O} & \mathbb{O} & 0 & i\eta_{s}%
\end{array}
\right)  \left(
\begin{array}
[c]{c}%
\left\vert v_{1}^{\ast}\right\rangle \\
\left\vert v_{1}\right\rangle \\
\vdots\\
\left\vert v_{s}^{\ast}\right\rangle \\
\left\vert v_{s}\right\rangle
\end{array}
\right)
\end{equation}


The bases $\left(  \left\vert v_{1}\right\rangle \left\vert v_{1}^{\ast
}\right\rangle \cdots\left\vert v_{s}\right\rangle \left\vert v_{s}^{\ast
}\right\rangle \right)  $ and $\left(  \left\vert s_{1}\right\rangle
\left\vert s_{1+s}\right\rangle \cdots\left\vert s_{s}\right\rangle \left\vert
s_{s+s}\right\rangle \right)  $ can be viewed as transformations of a fixed
\emph{real }starting basis $\left(  \left\vert e_{1}\right\rangle
\cdots\left\vert e_{2s}\right\rangle \right)  $ i.e.%
\begin{align}
\left(  \left\vert s_{1}\right\rangle \left\vert s_{1+s}\right\rangle
\cdots\left\vert s_{s}\right\rangle \left\vert s_{s+s}\right\rangle \right)
& =\left(  \left\vert e_{1}\right\rangle \cdots\left\vert e_{2s}\right\rangle
\right)  \mathbb{S}\nonumber\\
\left(  \left\vert v_{1}\right\rangle \left\vert v_{1}^{\ast}\right\rangle
\cdots\left\vert v_{s}\right\rangle \left\vert v_{s}^{\ast}\right\rangle
\right)   & =\left(  \left\vert e_{1}\right\rangle \cdots\left\vert
e_{2s}\right\rangle \right)  \mathbb{V}%
\end{align}
where $\mathbb{S}$ is a real orthogonal matrix and $\mathbb{V}$ is unitary
matrix. If we are solely working with matrices we can set
\begin{equation}
\left\vert e_{j}\right\rangle =\left(
\begin{array}
[c]{c}%
0\\
\vdots\\
0\\
1\\
0\\
\vdots\\
0
\end{array}
\right)  \leftarrow\text{j'th position}%
\end{equation}
and $\mathbb{S}=\left(  \left\vert s_{1}\right\rangle \left\vert
s_{1+s}\right\rangle \cdots\left\vert s_{s}\right\rangle \left\vert
s_{s+s}\right\rangle \right)  $ and $\mathbb{V}=$ $\left(  \left\vert
v_{1}\right\rangle \left\vert v_{1}^{\ast}\right\rangle \cdots\left\vert
v_{s}\right\rangle \left\vert v_{s}^{\ast}\right\rangle \right)  $ where
$\left\vert s_{j}\right\rangle $ and $\left\vert v_{j}\right\rangle $ are
numerical column vectors, while if we are considering abstract operators we
can consider $\left\{  \left\vert e_{j}\right\rangle \right\}  $ to be a fixed
real given basis of $\mathcal{H}$.

\section{Symmetry}

There is a 1-1 correspondence between antisymmetric operators acting on a
Hilbert space $\mathcal{H}$ (i.e. mixed contra and covariant tensors of rank
$\left(  1,1\right)  $) and contravariant tensors of rank $\left(  0,2\right)
$ and also covariant tensors of rank $\left(  2,0\right)  $. The first
correspondence is given by
\begin{equation}
\mathcal{A}=\sum_{1\leq j\leq s}\eta_{j}\left\{  \left\vert s_{j}\right\rangle
\left\langle s_{j+s}\right\vert -\left\vert s_{_{j+s}}\right\rangle
\left\langle s_{j}\right\vert \right\}  \leftrightarrow2\sum_{1\leq j\leq
s}\eta_{j}\left\vert s_{j}s_{j+s}\right\rangle \label{contra}%
\end{equation}
where $\left\vert s_{j}s_{j+s}\right\rangle $ is the normalized antisymmetric
tensor product of the \emph{real }basis elements $\left\{  \left\vert
s_{j}\right\rangle ,\left\vert s_{_{j+s}}\right\rangle \right\}  $ and
$\eta_{j}=\eta_{jj+s}$. The expansions of $\mathcal{A}$ are not unique and we
have from the preceeding that%
\begin{align}
\mathcal{A}  & =\sum_{\substack{1\leq j\leq s \\1\leq j^{\prime}\leq2s \\1\leq
j^{\prime\prime}\leq2s }}\eta_{j}\left\{  \left\vert e_{j^{\prime}%
}S_{j^{\prime}2j-1}\right\rangle \left\langle e_{j^{\prime\prime}}%
S_{j^{\prime\prime}2j}\right\vert -\left\vert e_{j^{\prime}}S_{j^{\prime}%
2j}\right\rangle \left\langle e_{j^{\prime\prime}}S_{j^{\prime\prime}%
2j-1}\right\vert \right\} \nonumber\\
& =\sum_{\substack{1\leq j\leq s \\1\leq j^{\prime}\leq2s \\1\leq
j^{\prime\prime}\leq2s }}\left\{  \eta_{j}\left\vert e_{j^{\prime}%
}S_{j^{\prime}2j-1}\right\rangle \left\langle e_{j^{\prime\prime}}%
S_{j^{\prime\prime}2j}\right\vert +\eta_{j+s}\left\vert e_{j^{\prime}%
}S_{j^{\prime}2j}\right\rangle \left\langle e_{j^{\prime\prime}}%
S_{j^{\prime\prime}2j-1}\right\vert \right\} \nonumber\\
& =\sum_{\substack{1\leq j\leq2s \\1\leq j^{\prime}\leq2s \\1\leq
j^{\prime\prime}\leq2s }}S_{j^{\prime}j}\eta_{j}S_{j^{\prime\prime}%
j}\left\vert e_{j^{\prime}}\right\rangle \left\langle e_{j^{\prime\prime}%
}\right\vert \nonumber\\
& =\sum_{\substack{1\leq j^{\prime}\leq2s \\1\leq j^{\prime\prime}\leq2s
}}\left(  \mathbb{S}\mathbf{\eta}\mathbb{S}^{t}\right)  _{_{j^{\prime
}j^{\prime\prime}}}\left\vert e_{j^{\prime}}\right\rangle \left\langle
e_{j^{\prime\prime}}\right\vert
\end{align}
and
\begin{align}
\mathcal{A}  & =2\sum_{1\leq j\leq s}\eta_{j}\left\vert s_{j}s_{j+s}%
\right\rangle =\sum_{1\leq j\leq s}\left\{  \eta_{jj+s}\left\vert s_{j}%
s_{j+s}\right\rangle +\eta_{j+sj}\left\vert s_{j+s}s_{j}\right\rangle \right\}
\nonumber\\
& =\sum_{\substack{1\leq j^{\prime}\leq2s \\1\leq j^{\prime\prime}\leq2s
}}\left(  \mathbb{S}\mathbf{\eta}\mathbb{S}^{t}\right)  _{_{j^{\prime
}j^{\prime\prime}}}\left\vert e_{j^{\prime}}e_{j^{\prime\prime}}\right\rangle
=\sum_{\substack{1\leq j^{\prime}\leq2s \\1\leq j^{\prime\prime}\leq2s
}}\mathbb{A}_{_{j^{\prime}j^{\prime\prime}}}\left\vert e_{j^{\prime}%
}e_{j^{\prime\prime}}\right\rangle
\end{align}
Eq. $\left(  \ref{contra}\right)  $ shows that
\begin{equation}
\mathcal{A}=2\sum_{1\leq j\leq s}\eta_{j}\left\vert Os_{j}Os_{j+s}%
\right\rangle =2\sum_{1\leq j\leq s}\eta_{j}\left\vert s_{j}s_{j+s}%
\right\rangle
\end{equation}
if
\begin{equation}
O\left(  \left\vert s_{j}\right\rangle ,\left\vert s_{j+s}\right\rangle
\right)  =\left(  \left\vert s_{j}\right\rangle ,\left\vert s_{j+s}%
\right\rangle \right)  \mathbb{O}_{\left(  j\right)  }%
\end{equation}
and
\begin{equation}
\det\mathbb{O}_{\left(  j\right)  }=1
\end{equation}
i.e. $O\in SO(2,\mathbb{R)}$ and $O\mathcal{A}O^{t}=\mathcal{A}.$ If we have
some degeneracies i.e. $\eta_{j_{1}}=\cdots=\eta_{j_{d}}$ then $O\in
SO(2d,\mathbb{R)}$ and $O\mathcal{A}O^{t}=\mathcal{A}$ also produce different
but equivalent expansions of $\mathcal{A}$ in terms of a real basis.

We can also observe that
\begin{equation}
\mathcal{A}=2\sum_{1\leq j\leq s}\eta_{j}\left\vert Us_{j}Us_{j+s}%
\right\rangle =2\sum_{1\leq j\leq s}\eta_{j}\left\vert s_{j}s_{j+s}%
\right\rangle
\end{equation}
if
\begin{equation}
U\left(  \left\vert s_{j}\right\rangle ,\left\vert s_{j+s}\right\rangle
\right)  =\left(  \left\vert s_{j}\right\rangle ,\left\vert s_{j+s}%
\right\rangle \right)  \mathbb{U}_{\left(  j\right)  }%
\end{equation}
and
\begin{equation}
\det\mathbb{U}_{\left(  j\right)  }=1
\end{equation}
i.e. $U\in SU(2,\mathbb{C)}$ and $U\mathcal{A}U^{t}=\mathcal{A}.$ If we have
some degeneracies i.e. $\eta_{j_{1}}=\cdots=\eta_{j_{d}}$ then $U\in
SU(2d,\mathbb{C)}$ and $U\mathcal{A}U^{t}=\mathcal{A}$ also produce different
but equivalent expansions of $\mathcal{A}$ in terms of a complex basis.

We can further observe more generally that
\begin{equation}
\mathcal{A}=2\sum_{1\leq j\leq s}\eta_{j}\left(  \det\mathbb{U}_{\left(
j\right)  }\right)  ^{-1}\left\vert Us_{j}Us_{j+s}\right\rangle =2\sum_{1\leq
j\leq s}\eta_{j}\left\vert s_{j}s_{j+s}\right\rangle
\end{equation}
if
\begin{equation}
U\left(  \left\vert s_{j}\right\rangle ,\left\vert s_{j+s}\right\rangle
\right)  =\left(  \left\vert s_{j}\right\rangle ,\left\vert s_{j+s}%
\right\rangle \right)  \mathbb{U}_{\left(  j\right)  }%
\end{equation}
and
\begin{equation}
\det\mathbb{U}_{\left(  j\right)  }=e^{i\sigma}%
\end{equation}
A special case of this is
\begin{equation}
\mathcal{A}=2\sum_{1\leq j\leq s}\eta_{j}\left(  \det\mathbb{U}_{\left(
j\right)  }\right)  ^{-1}\left\vert Us_{j}Us_{j+s}\right\rangle =2\sum_{1\leq
j\leq s}\left\vert \eta_{j}\right\vert \left\vert Us_{j}Us_{j+s}\right\rangle
\end{equation}
where
\begin{equation}
\eta_{j}=\left\vert \eta_{j}\right\vert e^{i\sigma}%
\end{equation}
Thus one can conclude that the expansion of $\mathcal{A}$ in terms of a
canonical complex basis \emph{does not }lead to \emph{unique }canonical
coefficients!! In this case canonical coefficients are only uniquely defined
up to a complex phase factor and the value of the canonical coefficients
depends on the canonical basis used.

Expansions in terms of canonical real bases as $\det\mathbb{O}_{\left(
j\right)  }=-1$ (the only other possible value for $\det\mathbb{O}_{\left(
j\right)  })$ leads to%
\begin{equation}
\mathcal{A}=2\sum_{1\leq j\leq s}\eta_{j}\left(  -1\right)  \left\vert
Os_{j}Os_{j+s}\right\rangle =2\sum_{1\leq j\leq s}\eta_{j}\left\vert
Os_{j+s}Os_{j}\right\rangle
\end{equation}
which corresponds to $\eta_{j}\rightarrow-\eta_{j}$. Thus canonical
coefficients in this case are uniquely defined up to sign. The eigenvalues of
$\mathcal{A}$ are however \emph{uniquely} defined as $\left\{  \pm i\eta
_{j};1\leq j\leq s\right\}  $.

\section{Complex Conjugation\label{star}}

A complex conjugation map, $K$, is determined by any representative,$\left\{
\left\vert e_{j}\right\rangle \right\}  $, of an equivalence class of bases
$\left\{  \left\vert \mathbf{e}_{k}\right\rangle \right\}  $ of $\mathcal{H}$
by
\begin{equation}
K\left\vert v\right\rangle =%
%TCIMACRO{\dsum }%
%BeginExpansion
{\displaystyle\sum}
%EndExpansion
\left(  x_{j}-iy_{j}\right)  \left\vert e_{j}\right\rangle =%
%TCIMACRO{\dsum }%
%BeginExpansion
{\displaystyle\sum}
%EndExpansion
\eta_{j}^{\ast}\left\vert e_{j}\right\rangle
\end{equation}
where
\begin{equation}
\left\vert v\right\rangle =%
%TCIMACRO{\dsum }%
%BeginExpansion
{\displaystyle\sum}
%EndExpansion
\eta_{j}\left\vert e_{j}\right\rangle
\end{equation}
Any vector that satisfies
\begin{equation}
K\left\vert w\right\rangle =\left\vert w\right\rangle
\end{equation}
is called real wrt $K$, in particular the basis elements $\left\{  \left\vert
e_{j}\right\rangle \right\}  $ are real vectors.

The complex conjugate of an operator $\mathcal{B}$ is defined by
\begin{equation}
\mathcal{B}^{\ast}=K\mathcal{B}K
\end{equation}
leading to
\begin{equation}
\mathcal{B}^{\ast}=K\mathcal{B}K\left(  \left\vert \mathbf{e}\right\rangle
\right)  =K\mathcal{B}\left(  \left\vert \mathbf{e}\right\rangle \right)
=K\left(  \left\vert \mathbf{e}\right\rangle \right)  \mathbb{B}=\left(
\left\vert \mathbf{e}\right\rangle \right)  \mathbb{B}^{\ast}%
\end{equation}


The transpose $\mathcal{B}^{t}$ can then be defined in terms of $K$ and the
adjoint operation as
\begin{equation}
\mathcal{B}^{t}=\left(  \mathcal{B}^{\dagger}\right)  ^{\ast}=K\mathcal{B}%
^{\dagger}K
\end{equation}



\end{document}